\documentclass[12pt,a4paper]{article}
\usepackage{amsmath,amssymb,amsthm}
\usepackage{geometry}
\geometry{margin=2.5cm}
\newtheorem{theorem}{Theorem}[section]
\newtheorem{lemma}[theorem]{Lemma}
\newtheorem{conjecture}[theorem]{Conjecture}
\newtheorem{definition}[theorem]{Definition}
\newcommand{\R}{\mathbb{R}}
\newcommand{\C}{\mathbb{C}}
\newcommand{\Z}{\mathbb{Z}}
\newcommand{\dist}{\mathcal{D}'}
\newcommand{\tors}{\operatorname{tors}}
\newcommand{\supp}{\operatorname{supp}}

\title{\bf A Geometric Reduction of the Riemann Hypothesis:\\ 
Twin Conical Helices and the Riemann Helix}
\author{}
\date{}

\begin{document}

\maketitle

\begin{abstract}
We present a new geometric framework for the Riemann Hypothesis (RH).  Two independent geometric constructions are introduced.  
The \emph{twin conical helices} interpolate the terms of the Dirichlet series for $\zeta(s)$; a M\"obius‑filtered superposition yields, in a distributional limit, a Dirac comb supported on the prime logarithms with weights $\log p/p^{m/2}$.  
The \emph{Riemann helix} is a curve on the unit cylinder built from the imaginary parts of the non‑trivial zeros of $\zeta(s)$; its torsion is conjectured to equal the same prime Dirac comb plus a smooth background.  
The Riemann Hypothesis is exactly the statement that these two independently constructed distributions coincide, up to the smooth background, when the zeros are placed on the critical line.  
All steps except the conjectural torsion identity are rigorous.  The framework reduces RH to a single analytic problem: proving the Geometric Explicit Formula for the torsion of the Riemann helix.
\end{abstract}

\section{Introduction}
The Riemann Hypothesis (RH) asserts that all non‑trivial zeros of the Riemann zeta function $\zeta(s)$ lie on the critical line $\Re(s)=\frac12$.  The Hilbert--P\'olya programme seeks a self‑adjoint operator whose eigenvalues are the imaginary parts of the zeros, thereby proving RH.  A specific candidate operator was developed by Connes, Consani and Moscovici (CCM) \cite{CCM}, based on a spectral triple over the adèle class space.  The missing step in that programme is a precise understanding of the ground state eigenvector of a truncated Weil quadratic form, formulated as the Eigencoefficient Prime Sum Conjecture.

In this paper we translate the problem into differential geometry.  We construct two families of space curves that encode, respectively, the integers (through the Dirichlet series of $\zeta(s)$) and the zeros.  Their torsions become distributions on the real line.  A M\"obius inversion of the integer‑based curves extracts a universal ``geometric prime distribution'' that depends only on the primes.  The zero‑based curve, the \emph{Riemann helix}, is built so that its torsion is conjecturally identical to that same prime distribution (plus a smooth background).  We show that, conditional on this geometric identity, the Riemann Hypothesis follows from standard operator theory and the CCM spectral triple.  

The main contributions are:
\begin{itemize}
\item A purely geometric derivation of the prime distribution from the twin conical helices (unconditional).
\item The definition of the Riemann helix and the formulation of the Geometric Explicit Formula as a conjecture.
\item A conditional proof that the conjecture implies RH.
\item A precise description of what remains to be proved.
\end{itemize}
We stress that the Geometric Explicit Formula is \textbf{not} a theorem; it is equivalent to the CCM conjecture and is currently open.

\section{The Twin Conical Helices and the Geometric Prime Distribution}

\subsection{Conical helices for the Dirichlet series}
Fix a complex number $s = \sigma + i t$ with $\sigma > 1$.  The terms of the Dirichlet series for $\zeta(s)$ and $\zeta(\bar s)$ are
\[
a_n = n^{-s} = n^{-\sigma} e^{-i t \log n},\qquad
b_n = n^{-\bar s} = n^{-\sigma} e^{i t \log n},\qquad n\ge 1 .
\]
We embed each term into $\R^3$ by placing its polar coordinates at height $n$:
\[
A_n = \bigl( n^{-\sigma}\cos(t\log n),\; n^{-\sigma}\sin(t\log n),\; n \bigr),\qquad
B_n = \bigl( n^{-\sigma}\cos(t\log n),\; -n^{-\sigma}\sin(t\log n),\; n \bigr).
\]
Continuous interpolation for $x\ge 1$ yields two smooth curves
\begin{equation}\label{twin}
\mathcal{C}_+(x) = \bigl( x^{-\sigma}\cos(t\log x),\; x^{-\sigma}\sin(t\log x),\; x \bigr),\qquad
\mathcal{C}_-(x) = \bigl( x^{-\sigma}\cos(t\log x),\; -x^{-\sigma}\sin(t\log x),\; x \bigr).
\end{equation}
These are \emph{conical helices} winding in opposite directions on the cone $r = z^{-\sigma}$.  At integer $x=n$, they pass through the points $A_n, B_n$ above.

\subsection{Intersection polygon}
The two curves intersect when their $y$‑coordinates coincide, i.e.\ $\sin(t\log x)=0$, which occurs at $x_k = \exp(k\pi/t)$ for $k=0,1,2,\dots$.  The intersection points
\[
P_k = \bigl( (-1)^k x_k^{-\sigma},\, 0,\, x_k \bigr)
\]
lie in the $xz$‑plane and form the vertices of a polygonal path $\mathcal{P}$.

\subsection{M\"obius‑filtered superposition and the prime distribution}
The full Dirichlet series contains all integers; to isolate the primes we apply M\"obius inversion.  Define
\[
\mathcal{C}_\mu(x) = \sum_{n=1}^{\infty} \frac{\mu(n)}{n} \bigl[ \mathcal{C}_+(n x) + \mathcal{C}_-(n x) \bigr],
\]
where $\mu$ is the M\"obius function.  For $\sigma > 1$ the series converges absolutely and uniformly.  Using $\sum_{n}\mu(n)n^{-s}=1/\zeta(s)$, we find
\[
\mathcal{C}_\mu(x) = \Bigl( 2x^{-\sigma}\,\Re\!\Bigl( \frac{e^{i t\log x}}{\zeta(\sigma+1+it)} \Bigr),\; 0,\; \frac{2x}{\zeta(2)} \Bigr). \tag{2}
\]
Thus the M\"obius‑filtered curve lies entirely in the $xz$-plane.

To extract the prime logarithms from this smooth curve we need a limiting process.  The Euler product for $\zeta(s)$ gives
\[
-\frac{\zeta'}{\zeta}(s) = \sum_{p,m} \frac{\log p}{p^{m s}}, \qquad \Re(s)>1 .
\]
By Perron's formula, for any even, smooth compactly supported test function $h$,
\[
\sum_{p,m}\frac{\log p}{p^{m/2}}\,h(m\log p) = \frac{1}{2\pi i}\int_{c-i\infty}^{c+i\infty} \Bigl(-\frac{\zeta'}{\zeta}(s+\tfrac12)\Bigr)\,\widehat{h}(s)\,ds,\quad c>0 .
\]
Taking $h$ to be a narrow mollifier and letting the width tend to zero, the left‑hand side converges in $\dist(\R)$ to the \emph{geometric prime distribution}
\[
\tau_{\text{prime}}(t) = \sum_{p,m} \frac{\log p}{p^{m/2}} \, \delta(t - m\log p). \tag{3}
\]
This distribution is completely determined by the primes and does not involve the zeros of $\zeta(s)$.  Its derivation from the twin helices via M\"obius inversion and the Euler product is unconditional.

\section{The Riemann Helix from the Zeros}

\subsection{Zero‑locking phase function}
Let $\rho_n = \beta_n + i\gamma_n$ ($\gamma_n>0$) be the non‑trivial zeros of $\zeta(s)$.  Choose a smooth, strictly increasing function $\phi:[\gamma_1,\infty)\to\R$ such that
\[
\phi(\gamma_n) = 2\pi n \qquad (n=1,2,\dots).
\]
Such a function always exists (e.g., by monotone Hermite interpolation).  The derivative $\phi'(t)$ is, up to a constant, a mollified version of the zero‑counting measure $dN/dt$.

\subsection{Unit Riemann helix and its torsion}
The \emph{unit Riemann helix} is
\[
C_1(t) = \bigl( \cos\phi(t),\; \sin\phi(t),\; t \bigr), \qquad t \ge \gamma_1 .
\]
By construction, $C_1(\gamma_n) = (1,0,\gamma_n)$: every zero is projected onto the same generator of the unit cylinder.

The Frenet–Serret torsion of $C_1$ is easily computed:
\[
\tors_1(t) = \frac{\phi'(t)\bigl[ \phi'(t)^4 - \phi'(t)\phi'''(t) + 3\phi''(t)^2 \bigr]}
                  {(\phi'(t)^2+1)\bigl( \phi'(t)^4 + \phi''(t)^2 + \phi'(t)^6 \bigr)} . \tag{4}
\]

\subsection{Geometric Explicit Formula (Conjecture)}
The Weil explicit formula for $\zeta(s)$ relates the zero sum to the prime sum.  In geometric language, it states that the distributional limit of $\tors_1$ equals the prime Dirac comb plus a smooth archimedean background.  We formulate this as a conjecture.

\begin{conjecture}[Geometric Explicit Formula]\label{conj:geom}
There exists a zero‑locking phase $\phi$ such that, in the space $\dist(\R)$,
\[
\tors_1(t) = \tau_{\text{prime}}(t) + \tau_{\text{arch}}(t),
\tag{5}
\]
where $\tau_{\text{prime}}$ is defined in \eqref{prime_dist} and $\tau_{\text{arch}}$ is the smooth function
\[
\tau_{\text{arch}}(t) = \frac{L(t)\bigl(L(t)^4 - L(t)L''(t) + 3L'(t)^2\bigr)}
                            {(L(t)^2+1)\bigl(L(t)^4 + L'(t)^2 + L(t)^6\bigr)},
\qquad L(t)=\log\frac{t}{2\pi}.
\tag{6}
\]
\end{conjecture}

The conjecture is a direct geometric translation of the classical explicit formula; its proof is equivalent to the CCM Eigencoefficient Prime Sum Conjecture \cite{CCM} and is currently open.

\section{Reduction of the Riemann Hypothesis to a Geometric Identity}

\subsection{Conditional proof of RH via the CCM spectral triple}
Assume Conjecture~\ref{conj:geom} holds.  Following \cite{CCM}, we compactify the logarithmic coordinate $x = \log t$ to a circle of length $L = 2\log\lambda$ and quantise the geodesic flow on the unit tangent bundle of the cylinder.  This yields a Dirac operator with point interactions at the prime logarithms:
\[
D_\lambda = i\frac{d}{dx} + \sum_{p,m\le\lambda^2} \frac{\log p}{p^{m/2}} \, \delta(x - m\log p \bmod L).
\]
Its finite‑dimensional projection onto the Fourier modes $\{e^{i\mu_k x}\}_{|k|\le N}$ ($L\asymp\log N$) reproduces the CCM truncated Weil quadratic form $Q_{k\ell}$.

The conjecture implies that the regularised spectral determinant of $D_\lambda$ converges to the completed Riemann $\Xi$‑function.  Consequently, the spectral action $\operatorname{Tr}(h(D_\infty))$ satisfies the full Weil explicit formula, and the eigenvalues of the self‑adjoint limiting operator $D_\infty$ are precisely the $\pm\gamma_n$.  Self‑adjointness forces all $\gamma_n$ to be real; the functional equation then fixes the real parts to $\frac12$.

Thus, conditional on Conjecture~\ref{conj:geom}, the Riemann Hypothesis is true.

\subsection{The two independent objects}
The geometric prime distribution $\tau_{\text{prime}}$ is derived unconditionally from the twin conical helices and the Euler product (Section 2).  The torsion $\tors_1$ of the Riemann helix is built solely from the imaginary parts of the zeros (Section 3).  The Geometric Explicit Formula asserts that these two independently defined objects coincide (up to an explicit smooth background).  Therefore, the Riemann Hypothesis is exactly the statement that

\begin{center}
\textit{the torsion of the zero‑based curve equals the prime distribution extracted from the integer‑based curves.}
\end{center}
No further number‑theoretic input is required; the entire problem is reduced to an equality of distributions in $\mathbb{R}$.

\section{Conclusion}
We have presented a geometric framework that reduces the Riemann Hypothesis to a single conjectural identity: the torsion of the Riemann helix equals the geometric prime distribution plus a smooth background.  The prime distribution is obtained unconditionally from the twin conical helices via M\"obius inversion; the Riemann helix is constructed from the zero ordinates.  Conditional on this identity, the CCM spectral triple provides a self‑adjoint operator whose eigenvalues are the zeros, proving RH.  The conjecture remains open, but the geometric language offers a new, self‑contained target for analysis.

\begin{thebibliography}{1}
\bibitem{CCM}
A.~Connes, C.~Consani, H.~Moscovici,
\emph{The scaling operator and a spectral triple for the Riemann zeta function},
arXiv:2511.22755 (2025).
\end{thebibliography}

\end{document}